\documentclass{article}
\usepackage{amsmath,amssymb,hyperref}
\title{When a real retry opportunity persists across other work}
\date{26 November 2026}
\begin{document}
\maketitle
\paragraph{Claim and stakes.}
A civil inspection report is owed at a receiving application. The
fact that its pending send repeatedly becomes eligible does not make
the send continuously eligible. In the pinned simulator radio path,
another valid transmission can remove eligibility before the next
opportunity to attempt the report. An installable budget check can
prevent that particular loss without stopping disjoint traffic. It
does not force the consumer to act, preserve eligibility under a new
route failure, or certify deadline service without a scheduling bound.

\paragraph{What is native.}
Replay selected, byte-pinned functions at osrf/subt revision
\texttt{083c6a8acb0c25aa1a35e34517e628cd3b52ae77}, using the earlier
broker harness and its local fixtures. Native broker dispatch calls
radio accounting then the application callback. A transmission debits
the sending node before final path and receive acceptance; a successful
receive also debits the actual physical recipient. At time \(t\), ledger
rows strictly earlier than \(t-1\) are removed by the respective
accounting code; a row at equality still counts. Use integer byte
caps with fixed homogeneous configuration, stable addresses and
endpoint bindings. The duties, gates, service opportunities and
loss schedule are fixture inputs, not discovered infrastructure
features. An actual callback is the terminal effect of this model,
not physical arrival, custody transfer, or a human receiving a report.

Fix a named duty of \(s\) bytes, eligible at a cut if a dispatch there
can make its actual callback under the installed mechanism. Treat E
as a proof predicate on complete state, not a request for a ghost
oracle by the controller. Evidence of enabledness comes from native
dispatch on a separately reconstructed branch with the identical prior
history, checked against the snapshot ledger inequality. Never test
it by spending capacity in the live branch. T is occurrence of that
same duty's terminal callback, not queue presence, poll, sent datagram,
or transmission of another report.

\paragraph{Source witness and enforceable boundary.}
Use a two-second period. Open with the old obligation pending and
empty relevant rolling windows, when a real old dispatch would pay.
At \(+.5\) an optional report of \(r\) bytes dispatches, then there is
another old scheduling cut at \(+.75\). When \(s=r=3\) and cap \(c=5\),
either shared sender or shared physical receive resource can make
the late cut fail, despite the first cut permitting the old report.
Different application names can share the same physical receive row.
A lost optional transmission still consumes its shared sending row;
it does not consume a receive row in this selected erased branch.
An unrelated application uses disjoint rows and makes an actual
delivery in every period. Logical repetition follows because all
charged rows in this fixture are outside the rolling window before
the next opening, bindings and offered work are the same up to time
translation, and the app's unobserved counters do not feed the policy.
Four native repetitions are checked; extrapolation uses this fixed
repeatable model, not a long-run network measurement.

For \(c\) fixed, the reserve checks potential competing
increments \(\Delta_T,\Delta_R\) at the two relevant ledgers and
retains the old \(s\) bytes of headroom:
\[
  L_T+\Delta_T+s\le c,\qquad L_R+\Delta_R+s\le c.
\]
A guard at the actual consuming cut admits a competing effect only
if it satisfies both, in all admitted route branches. For this
family the rule is sharp for \emph{immediate persistence} when those
maximal debits can occur. Letting any witnessed over-budget shared
debit occur disables the corresponding path until enough rows expire.
No budget is pre-debited to fake a physical payment. Time without
new debits can only free capacity. Thus if the gate covers all relevant
intervening debit writes until actual use, the route, permission,
address and capacity qualification are stable, and old was already
eligible, any intervening non-use step preserves its eligibility.
The control is new outer fixture logic, not a guard known to exist
in stock SubT. A lease that expires exactly when the optional writer
can spend capacity fails in the replay; one that outlives that one
writer succeeds on these paths but is not a permanent invariant.
Changing the old route's physical availability disables the action
even with the budget held. A permanent finite-service resource
tradeoff remains: blocking an optional conflicting report costs
that report unless it can be rescheduled or redirected. A disjoint
optional destination and origin still pass. A direct reservation
scheduler with the same scopes, information and enforcement ties.

\paragraph{Weak and strong fairness are different warranties.}
On an infinite trace, take weak fairness as recurring
\(\neg E\lor T\), and strong fairness as recurring \(E\) implying
recurring \(T\). This is the single-action distinction in
\href{https://lamport.azurewebsites.net/tla/tutorial/session11-2.html}{Lamport's
fairness tutorial}, with \(E,T\) here attached to the simulator cuts.
The indefinitely unpaid period above passes weak fairness (each late
cut disables it) and violates strong fairness (each opening enables
it). A scheduler that attempts fresh optional work and retries the
old report only at the second cut is a witness in the model, not a
claim about a stock automated scheduler. Adding strong fairness as
a verifier assumption would delete that valid fixture history without
adding any scheduler, reserve, or permission to make the first cut
usable. The earlier
\href{https://github.com/eternal-november/eos-research-programmes/pull/617}{recurrence
checker} can encode weak fairness as an edge-justice label; strong
fairness is not equivalent for transient eligibility.

There is a sufficient replacement law. Suppose every non-\(T\)
transition preserves \(E\) once true. Strong always implies weak;
conversely if a weakly fair run visits \(E\) infinitely but only
finitely many \(T\), after the last \(T\) there is a first \(E\),
then persistence makes \(E\) hold forever, contradicting weak
fairness. The guard yields this premise on the qualified deployment
envelope, not weak fairness itself. No finite bound follows:
an arbitrary finite number of unpaid periods followed by a real
payment is consistent with any such eventual-service interpretation.
If the admitted controller is promised to try a persistently
enabled duty at the first late cut, the model pays at that cut.
The scheduling promise must be established independently.
The full negative test leaves the old duty unpaid but still admits a
paid continuation from each prefix; this is not a vacuity from
assuming the source transition graph stops at the problematic cut.

\paragraph{Verification and limits.}
The replay sweeps 12,992 authored world histories with four periods
and checks 103,936 actual candidate dispatches at proof-only forks
against a separate closed-form phase calculation. It distinguishes an
erased optional transfer with shared sender from an erased optional
transfer with only shared receiver. Lean checks the two logical
directions under their respective premises and the exact budget
predicate; the simulator and environment hypotheses remain in the
refinement argument. The result neither generalizes the \emph{strong}
fairness checker nor establishes operational superiority over a
stronger scheduler, unbounded stock simulator progress, field radio
behavior, a new inventory disposition, or independent mission acceptance.
\end{document}
